\clearpage
\VerdierPhysicalPageAnchor{0145}%
\begingroup
\small
\setlength{\parskip}{0pt}%
\setlength{\emergencystretch}{1.5em}%
\setlength{\abovedisplayskip}{3pt plus 1pt minus 1pt}%
\setlength{\belowdisplayskip}{3pt plus 1pt minus 1pt}%

\VerdierSemanticAnchor{ch2-2-3-2-free-object-continuation}
\noindent univers dont $C$ et $C'$ soient des éléments, en disant que $Y$
est libre à droite sur $X$ s'il existe un isomorphisme entre les foncteurs
$Z\longmapsto\operatorname{Hom}_C(Z,Y)$ et
$Z\longmapsto\operatorname{Hom}_{C'}(F(Z),X)$.

\VerdierSemanticAnchor{ch2-2-3-2-right-left-terminology}
Nous dirons qu'un objet $Y$ de $C$ est \emph{$F$-libre à droite} s'il existe
un objet $X$ de $C'$ tel que $Y$ soit libre à droite sur $X$. Nous dirons
qu'un objet $X$ de $C'$ est \emph{$F$-libérable à droite} s'il existe un
objet $Y$ de $C$ libre à droite sur $X$. On définit de façon analogue, par
passage aux catégories opposées, les notions d'objet libre ou libérable à
\emph{gauche}.

\VerdierSemanticAnchor{ch2-2-3-3}
\medskip
\begingroup\itshape
\noindent\textbf{Proposition 2.3.3.}\quad
Soient $\mathcal D$ une catégorie triangulée, $\mathcal B$ une
sous-catégorie triangulée pleine,
$Q\colon\mathcal D\longrightarrow\mathcal D/\mathcal B$ le foncteur de
passage au quotient (2.2.10), $S_{\mathcal D}(\mathcal B)$ le système
multiplicatif de $\mathcal D$ associé à $\mathcal B$ (2.1.8).

\VerdierSemanticAnchor{ch2-2-3-3-a}
\par\noindent a) Pour tout objet $Y$ de $\mathcal D$, les conditions
suivantes sont équivalentes :

\VerdierSemanticAnchor{ch2-2-3-3-a-i}
\par\noindent\hspace{1.2em}i) L'objet $Y$ est $Q$-libre à droite.

\VerdierSemanticAnchor{ch2-2-3-3-a-ii}
\par\noindent\hspace{1.2em}ii) Pour tout morphisme
$s\in S_{\mathcal D}(\mathcal B)$, $s\colon S\longrightarrow T$,
l'application :
\[
  \operatorname{Hom}_{\mathcal D}(s,Y)\colon
  \operatorname{Hom}_{\mathcal D}(T,Y)\longrightarrow
  \operatorname{Hom}_{\mathcal D}(S,Y)
\]
est une bijection.

\VerdierSemanticAnchor{ch2-2-3-3-a-iii}
\par\noindent\hspace{1.2em}iii) Pour tout objet $B$ de $\mathcal B$,
$\operatorname{Hom}_{\mathcal D}(B,Y)=0$.

\VerdierSemanticAnchor{ch2-2-3-3-a-iv}
\par\noindent\hspace{1.2em}iv) Tout morphisme
$Y\xrightarrow{s}Z$, avec $s\in S_{\mathcal D}(\mathcal B)$, admet une
rétraction.

\VerdierSemanticAnchor{ch2-2-3-3-a-v}
\par\noindent\hspace{1.2em}v) Pour tout objet $X$ de $\mathcal D$,
l'application :
\[
  \operatorname{Hom}_{\mathcal D}(X,Y)\longrightarrow
  \operatorname{Hom}_{\mathcal D/\mathcal B}(Q(X),Q(Y))
\]
est bijective.

\VerdierSemanticAnchor{ch2-2-3-3-b}
\par\noindent b) Soit $L_R(Q)$ la sous-catégorie pleine de $\mathcal D$
définie par les objets $Q$-libres à droite. La catégorie $L_R(Q)$ est
triangulée, strictement pleine, saturée.

\VerdierSemanticAnchor{ch2-2-3-3-c}
\par\noindent c) La restriction du foncteur de passage au quotient à
$L_R(Q)$ :
\[
  Q|_{L_R(Q)}\colon L_R(Q)\longrightarrow\mathcal D/\mathcal B
  ,
\]
est pleinement fidèle et injective sur les objets.

\VerdierSemanticAnchor{ch2-2-3-3-d}
\par\noindent d) Soit $\operatorname{Lib}_R(Q)\subset\mathcal D/\mathcal B$
la sous-catégorie pleine définie par les objets $Q$-libérables à droite de
$\mathcal D/\mathcal B$. La sous-catégorie $\operatorname{Lib}_R(Q)$ est
triangulée strictement pleine. Si les produits ou les sommes directes
dénombrables sont représentables dans $\mathcal D$, la catégorie
$\operatorname{Lib}_R(Q)$ est saturée.

\VerdierSemanticAnchor{ch2-2-3-3-e}
\par\noindent e) Le foncteur $Q$ envoie $L_R(Q)$ dans
$\operatorname{Lib}_R(Q)$. La catégorie $\operatorname{Lib}_R(Q)$ est la
clôture strictement pleine de l'image de $L_R(Q)$. Les flèches de liberté
sont des isomorphismes.
\endgroup
\endgroup
